\documentclass[conference]{IEEEtran}

% ---- Packages ----
\usepackage{cite}
\usepackage{amsmath,amssymb,amsfonts}
\usepackage{graphicx}
\usepackage{textcomp}
\usepackage{xcolor}
\usepackage{booktabs}
\usepackage{hyperref}
\usepackage{float}

\begin{document}

% ============================================================
%  TITLE & AUTHORS
% ============================================================
\title{Breast Cancer Diagnosis Using Classical Machine Learning:
       A Comparative Study with Ensemble Improvements}

\author{
  \IEEEauthorblockN{Your Full Name}
  \IEEEauthorblockA{
    \textit{Department of Computer Science} \\
    \textit{Your University Name} \\
    Your City, Your Country \\
    your.email@university.edu
  }
}

\maketitle


% ============================================================
%  ABSTRACT
% ============================================================
\begin{abstract}
Breast cancer remains one of the leading causes of cancer-related 
mortality worldwide, making early and accurate diagnosis critically 
important. This paper presents a comprehensive comparative study of 
classical machine learning methods applied to the UCI Wisconsin 
Breast Cancer Diagnostic (WDBC) dataset for binary classification 
of tumors as malignant or benign. Five baseline models are evaluated: 
Logistic Regression, Support Vector Machine (SVM), K-Nearest Neighbors 
(KNN), Decision Tree, and Naive Bayes. Building upon the baseline 
results, improvements are proposed through hyperparameter optimization 
via GridSearchCV, ensemble methods including Random Forest and Gradient 
Boosting, and a Soft Voting Ensemble combining the three best-performing 
models. Feature selection using SelectKBest reduced dimensionality from 
30 to 15 features. The proposed Voting Ensemble achieved an F1 score of 
[YOUR SCORE] and ROC-AUC of [YOUR SCORE], outperforming all baseline 
models. Results demonstrate that systematic hyperparameter tuning and 
ensemble strategies significantly improve diagnostic accuracy in 
classical ML pipelines.
\end{abstract}

\begin{IEEEkeywords}
breast cancer diagnosis, machine learning, SVM, ensemble methods, 
feature selection, classification, UCI WDBC dataset
\end{IEEEkeywords}


% ============================================================
%  SECTION 1 — INTRODUCTION
% ============================================================
\section{Introduction}

Breast cancer is one of the most prevalent and life-threatening 
diseases affecting millions of people globally. According to the 
World Health Organization, breast cancer accounts for approximately 
12\% of all new cancer diagnoses worldwide each year~\cite{who2023}. 
Early and accurate diagnosis is essential, as it significantly 
improves treatment outcomes and survival rates.

Machine learning (ML) has emerged as a powerful tool to assist 
medical professionals in diagnosing breast cancer from clinical 
measurements. Unlike deep learning approaches, classical ML methods 
offer interpretability and computational efficiency, making them 
well-suited for clinical decision support systems where transparency 
is critical.

This study uses the publicly available UCI Wisconsin Breast Cancer 
Diagnostic (WDBC) dataset~\cite{uci_wdbc}, which contains 569 
samples with 30 numerical features derived from digitized images 
of fine needle aspirate (FNA) of breast masses. Each sample is 
labeled as either malignant (212 samples) or benign (357 samples).

The primary contributions of this paper are:
\begin{itemize}
  \item A systematic evaluation of five classical ML classifiers 
        as baselines on the WDBC dataset.
  \item Proposed improvements through hyperparameter tuning, 
        ensemble methods, and feature selection.
  \item A noise robustness analysis evaluating model stability 
        under varying levels of Gaussian noise.
  \item Identification of the most discriminative features for 
        breast cancer classification.
\end{itemize}

The remainder of this paper is organized as follows. Section II 
reviews related work. Section III describes the methodology. 
Section IV presents results. Section V discusses findings, and 
Section VI concludes the paper.


% ============================================================
%  SECTION 2 — RELATED WORK
% ============================================================
\section{Related Work}

Several studies have applied classical ML methods to breast cancer 
diagnosis using the UCI WDBC dataset. Wolberg et al.~\cite{wolberg1995} 
originally introduced the dataset and demonstrated the utility of 
ML-based approaches for cytological diagnosis.

Akay~\cite{akay2009} proposed an SVM-based approach with feature 
selection achieving approximately 99.51\% accuracy on the WDBC 
dataset. Their work highlighted the importance of removing 
irrelevant features to improve SVM performance.

The PMLB benchmark study by Olson et al.~\cite{olson2017} 
evaluated a wide range of classifiers across many UCI datasets, 
establishing competitive baselines for ensemble methods such as 
Random Forest and Gradient Boosting on medical classification tasks.

Exarchos et al.~\cite{exarchos2012} demonstrated that combining 
multiple classifiers through ensemble strategies consistently 
outperforms individual models for clinical datasets, motivating 
the ensemble approach adopted in this study.

In contrast to these works, this paper focuses not only on 
achieving high accuracy but also on systematic improvement 
analysis, feature importance evaluation, and robustness testing 
under noisy conditions.


% ============================================================
%  SECTION 3 — METHODOLOGY
% ============================================================
\section{Methodology}

\subsection{Dataset}

The UCI Wisconsin Breast Cancer Diagnostic dataset~\cite{uci_wdbc} 
consists of 569 instances with 30 real-valued features computed 
from digitized FNA images. Features describe characteristics of 
cell nuclei including radius, texture, perimeter, area, smoothness, 
compactness, concavity, symmetry, and fractal dimension — each 
measured as mean, standard error, and worst value. The binary 
target variable indicates malignant (0) or benign (1) diagnosis.

\subsection{Preprocessing}

The dataset contained no missing values. The following preprocessing 
steps were applied:

\begin{enumerate}
  \item \textbf{Stratified Train/Test Split:} The dataset was 
    divided into 80\% training (455 samples) and 20\% testing 
    (114 samples) using stratified sampling to preserve class 
    distribution in both subsets.
  \item \textbf{Feature Scaling:} StandardScaler was applied to 
    normalize all features to zero mean and unit variance. 
    Critically, the scaler was fitted exclusively on the training 
    set and applied to the test set to prevent data leakage.
  \item \textbf{Feature Selection:} SelectKBest with ANOVA F-test 
    (\texttt{f\_classif}) was used to select the top 15 most 
    discriminative features from the original 30, reducing 
    dimensionality while retaining the most informative signals.
\end{enumerate}

\subsection{Baseline Models}

Five classical ML classifiers were trained using default 
scikit-learn parameters as baselines:

\begin{itemize}
  \item \textbf{Logistic Regression (LR):} A linear probabilistic 
    classifier well-suited for binary classification.
  \item \textbf{Support Vector Machine (SVM):} Uses an RBF kernel 
    to find an optimal separating hyperplane in a transformed 
    feature space.
  \item \textbf{K-Nearest Neighbors (KNN):} A non-parametric 
    classifier using k=5 nearest samples by Euclidean distance.
  \item \textbf{Decision Tree (DT):} A tree-based model that 
    recursively partitions the feature space.
  \item \textbf{Naive Bayes (NB):} A probabilistic classifier 
    based on Bayes' theorem with feature independence assumption.
\end{itemize}

\subsection{Proposed Improvements}

Four improvement strategies were applied:

\begin{enumerate}
  \item \textbf{Hyperparameter Tuning:} GridSearchCV with 5-fold 
    cross-validation optimized hyperparameters for LR 
    (C, solver, penalty) and SVM (C, gamma, kernel).
  \item \textbf{Random Forest:} An ensemble of decision trees 
    trained with bagging, tuned for n\_estimators, max\_depth, 
    and min\_samples\_split.
  \item \textbf{Gradient Boosting:} A sequential ensemble method 
    tuned for n\_estimators, learning\_rate, and max\_depth.
  \item \textbf{Soft Voting Ensemble:} Combined the Tuned LR, 
    Tuned SVM, and Random Forest using soft voting — averaging 
    predicted class probabilities to make final decisions.
\end{enumerate}

\subsection{Evaluation Metrics}

All models were evaluated using: Accuracy, Precision, Recall, 
F1 Score, ROC-AUC, and 5-fold Stratified Cross-Validation F1. 
In medical diagnosis, Recall (sensitivity) is particularly 
important as it measures the ability to correctly identify 
malignant cases — false negatives carry higher clinical risk 
than false positives.

\subsection{Noise Robustness}

To assess model stability, Gaussian noise at levels 
$\sigma \in \{0.0, 0.05, 0.1, 0.2, 0.3, 0.5\}$ was added 
to the test set and accuracy/F1 were recorded for the 
best-performing model.


% ============================================================
%  SECTION 4 — RESULTS
% ============================================================
\section{Results}

\subsection{Baseline Model Performance}

Table~\ref{tab:baseline} summarizes the performance of the five 
baseline classifiers. SVM and Logistic Regression achieved the 
highest F1 scores, consistent with the linear separability of 
the WDBC feature space after scaling.

% --- REPLACE VALUES WITH YOUR ACTUAL RESULTS ---
\begin{table}[H]
\centering
\caption{Baseline Model Performance}
\label{tab:baseline}
\begin{tabular}{lcccc}
\toprule
\textbf{Model} & \textbf{Acc} & \textbf{Prec} & 
\textbf{Recall} & \textbf{F1} \\
\midrule
Logistic Regression & 0.XXXX & 0.XXXX & 0.XXXX & 0.XXXX \\
SVM                 & 0.XXXX & 0.XXXX & 0.XXXX & 0.XXXX \\
KNN                 & 0.XXXX & 0.XXXX & 0.XXXX & 0.XXXX \\
Decision Tree       & 0.XXXX & 0.XXXX & 0.XXXX & 0.XXXX \\
Naive Bayes         & 0.XXXX & 0.XXXX & 0.XXXX & 0.XXXX \\
\bottomrule
\end{tabular}
\end{table}

\subsection{Improved Model Performance}

Table~\ref{tab:improved} presents the results after applying 
hyperparameter tuning and ensemble strategies. The Soft Voting 
Ensemble achieved the best overall performance across all metrics.

% --- REPLACE VALUES WITH YOUR ACTUAL RESULTS ---
\begin{table}[H]
\centering
\caption{Improved Model Performance}
\label{tab:improved}
\begin{tabular}{lcccc}
\toprule
\textbf{Model} & \textbf{Acc} & \textbf{Prec} & 
\textbf{Recall} & \textbf{F1} \\
\midrule
Tuned LR          & 0.XXXX & 0.XXXX & 0.XXXX & 0.XXXX \\
Tuned SVM         & 0.XXXX & 0.XXXX & 0.XXXX & 0.XXXX \\
Random Forest     & 0.XXXX & 0.XXXX & 0.XXXX & 0.XXXX \\
Gradient Boosting & 0.XXXX & 0.XXXX & 0.XXXX & 0.XXXX \\
Voting Ensemble   & 0.XXXX & 0.XXXX & 0.XXXX & 0.XXXX \\
\bottomrule
\end{tabular}
\end{table}

\subsection{ROC Curves}

Fig.~\ref{fig:roc} illustrates the ROC curves for all improved 
models. The Voting Ensemble and Tuned SVM achieved the highest 
AUC values, indicating strong discriminative ability between 
malignant and benign classes.

\begin{figure}[H]
  \centering
  \includegraphics[width=\columnwidth]{../final_roc_curves.png}
  \caption{ROC curves for improved models.}
  \label{fig:roc}
\end{figure}

\subsection{Confusion Matrix}

Fig.~\ref{fig:cm} shows the confusion matrix for the best 
performing model (Voting Ensemble). The model correctly 
classified [X] out of [X] malignant cases, yielding a 
false negative rate of [X]\%.

\begin{figure}[H]
  \centering
  \includegraphics[width=0.75\columnwidth]
                  {../best_model_confusion_matrix.png}
  \caption{Confusion matrix — Voting Ensemble.}
  \label{fig:cm}
\end{figure}

\subsection{Feature Importance}

Fig.~\ref{fig:fi} presents feature importances derived from 
the Random Forest model. [YOUR TOP FEATURE] emerged as the 
most discriminative feature, followed by [2ND FEATURE] and 
[3RD FEATURE], consistent with clinical understanding that 
tumor size and shape irregularity are key malignancy indicators.

\begin{figure}[H]
  \centering
  \includegraphics[width=\columnwidth]{../feature_importance.png}
  \caption{Feature importance scores from Random Forest.}
  \label{fig:fi}
\end{figure}

\subsection{Cross-Validation}

Fig.~\ref{fig:cv} shows the 5-fold stratified cross-validation 
F1 distributions. Improved models exhibit higher median scores 
and lower variance, indicating more stable generalization.

\begin{figure}[H]
  \centering
  \includegraphics[width=\columnwidth]{../cv_boxplot.png}
  \caption{5-fold CV F1 score distributions for all models.}
  \label{fig:cv}
\end{figure}

\subsection{Noise Robustness}

The Tuned SVM maintained [X]\% accuracy at noise level 
$\sigma=0.1$, degrading to [X]\% at $\sigma=0.5$, demonstrating 
reasonable robustness for clinical applicability where measurement 
noise is expected.


% ============================================================
%  SECTION 5 — DISCUSSION
% ============================================================
\section{Discussion}

\subsection{Why SVM and Logistic Regression Perform Best}

The strong baseline performance of SVM and Logistic Regression 
suggests the WDBC feature space is largely linearly separable 
after StandardScaler normalization. Distance-based and 
margin-based models particularly benefit from proper feature 
scaling, which is why these models outperformed the Decision 
Tree — a scale-invariant model.

\subsection{Impact of Hyperparameter Tuning}

GridSearchCV improved F1 by [X]\% for LR and [X]\% for SVM 
over their default configurations. This confirms that default 
hyperparameters are rarely optimal and systematic search is 
a worthwhile investment even for well-established models.

\subsection{Ensemble Methods}

Random Forest reduced overfitting compared to the single 
Decision Tree baseline by aggregating predictions across 
[200] trees with bootstrap sampling. The Soft Voting Ensemble 
further improved performance by leveraging the complementary 
strengths of LR, SVM, and Random Forest — each model 
contributes diverse decision boundaries.

\subsection{Clinical Implications of Recall}

In breast cancer diagnosis, a false negative (predicting benign 
when actually malignant) carries far greater clinical risk than 
a false positive. The Voting Ensemble achieved a Recall of 
[YOUR RECALL], meaning it correctly identified [X]\% of all 
malignant tumors — a critical metric for medical deployment.

\subsection{Limitations and Future Work}

Several limitations should be noted. First, the dataset size 
(569 samples) is relatively small; performance on larger 
clinical datasets remains to be validated. Second, deep 
learning methods were excluded by design — future work could 
compare classical ML against shallow neural networks. Third, 
class imbalance (357 benign vs 212 malignant) was addressed 
through stratified splits but not through oversampling 
techniques such as SMOTE, which could further improve Recall.


% ============================================================
%  SECTION 6 — CONCLUSION
% ============================================================
\section{Conclusion}

This paper presented a systematic study of classical machine 
learning methods for breast cancer diagnosis using the UCI 
WDBC dataset. Five baseline models were evaluated, with SVM 
and Logistic Regression achieving the strongest baseline 
performance. Four improvement strategies — hyperparameter 
tuning, Random Forest, Gradient Boosting, and a Soft Voting 
Ensemble — were applied and evaluated.

The proposed Voting Ensemble achieved the best overall 
performance with an F1 score of [YOUR SCORE], accuracy of 
[YOUR SCORE], and ROC-AUC of [YOUR SCORE], representing a 
meaningful improvement over the best baseline. Feature 
importance analysis identified [TOP FEATURE] as the most 
discriminative variable, consistent with clinical knowledge.

These results demonstrate that well-tuned classical ML 
pipelines with proper preprocessing and ensemble strategies 
can achieve clinically competitive diagnostic accuracy without 
the computational overhead of deep learning. Future work will 
explore SMOTE-based oversampling, larger clinical datasets, 
and explainability techniques such as SHAP values for 
clinical deployment.


% ============================================================
%  REFERENCES
% ============================================================
\begin{thebibliography}{00}

\bibitem{who2023}
World Health Organization, ``Breast Cancer,'' 
\textit{WHO Fact Sheets}, 2023. [Online]. Available: 
https://www.who.int/news-room/fact-sheets/detail/breast-cancer

\bibitem{uci_wdbc}
W. H. Wolberg, W. N. Street, and O. L. Mangasarian, 
``Breast Cancer Wisconsin (Diagnostic) Data Set,'' 
\textit{UCI Machine Learning Repository}, 1995.

\bibitem{wolberg1995}
W. H. Wolberg, W. N. Street, D. M. Heisey, and 
O. L. Mangasarian, ``Computerized breast cancer diagnosis 
and prognosis from fine-needle aspirates,'' 
\textit{Archives of Surgery}, vol. 130, no. 5, 
pp. 511--516, 1995.

\bibitem{akay2009}
M. F. Akay, ``Support vector machines combined with 
feature selection for breast cancer diagnosis,'' 
\textit{Expert Systems with Applications}, vol. 36, 
no. 2, pp. 3240--3247, 2009.

\bibitem{olson2017}
R. S. Olson, W. La Cava, P. Orzechowski, R. J. Urbanowicz, 
and J. H. Moore, ``PMLB: A large benchmark suite for machine 
learning evaluation and comparison,'' 
\textit{BioData Mining}, vol. 10, no. 36, 2017.

\bibitem{exarchos2012}
K. P. Exarchos, Y. Goletsis, and D. I. Fotiadis, 
``Multiparametric decision support system for the prediction 
of oral cancer reoccurrence,'' 
\textit{IEEE Transactions on Information Technology 
in Biomedicine}, vol. 16, no. 6, pp. 1127--1134, 2012.

\bibitem{breiman2001}
L. Breiman, ``Random forests,'' 
\textit{Machine Learning}, vol. 45, pp. 5--32, 2001.

\end{thebibliography}

\end{document}